%%%%%%%%%%%%%%%%%%%%%%%%%%%%%%%%%%%%%%%%%%%%%%%%%%%%%%%%%%%%%%%%%%%%%%%%%%%
\chapter{Review of Related Literature}
%%%%%%%%%%%%%%%%%%%%%%%%%%%%%%%%%%%%%%%%%%%%%%%%%%%%%%%%%%%%%%%%%%%%%%%%%%%

%%%%%%%%%%%%% SECTION 1 %%%%%%%%%%%%%
\section{Cognitive Load Theory}
Cognitive load theory (CLT) studies the strain induced by tasks on human information processing \cite{hart2006, sweller2019architecture}. This theory is based on the assumption that working memory is limited \cite{sweller2010element} and that cognitive overload can reduce performance \cite{howie2023} and learning \cite{sweller2019architecture, klepsch2022}. Experimental data have shown the empirical merits of this theory for understanding human processing of information \cite{ahmad2023} and are applied in the field of education \cite{elglaly2024} and technology \cite{klepsch2022}. Applying this theory to modern technology has shown effectiveness \cite{klepsch2022} and has the potential to create systems that better suit human cognitive architecture.

\subsection{Types of Cognitive Load}
Modern understanding of cognitive load is a three-part model, consisting of intrinsic, extraneous, and  germane cognitive load \cite{sweller1988cognitive, sweller2019architecture, klepsch2022}. Intrinsic cognitive load is load induced by the complexity of information being processed \cite{sweller2019architecture}; this kind of cognitive load is affected both by the information being processed and the expertise of the person receiving the information, since both these factors affect how complex the information being processed is perceived \cite{sweller2019architecture}. A domain expert will have an easier time understanding material in his/her field and thus will have less intrinsic cognitive load. Extraneous cognitive load is another component and is the load induced by the type of medium used to present information \cite{sweller2019architecture}, and other outside factors such noise, comfort and distractions \cite{howie2023}. Different mediums have different effects on cognitive load and can be observed in the split attention effect \cite{sweller2019architecture, thees2020}, and modality effect \cite{sweller2019architecture} of cognitive load research. Germane load is the load associated with schema acquisition or learning \cite{sweller1988cognitive, sweller2019architecture}. Working memory allocated to learning is considered the portion called germane load \cite{sweller2010element}, and is present in cases of moderate cognitive load. Cognitive overload does not allow germane load to occur \cite{sweller2019architecture}, thus little or no schema acquisition happens during cognitive overload.

\subsection{Factors Affecting Cognitive Load}
Experts and novices are differentiated by the amount of schema they have on a particular area of knowledge, not on the amount of working memory they have \cite{hart2006}. The more schema or background knowledge a person has, the more capable they are of solving novel or complex problems in the field by relying on existing strategies for similar problems \cite{hart2006}. Access to long-term memory is infinite, and thus the reliance of existing strategies in long- term memory does not induce high amounts of cognitive resources \cite{hart2006, sweller2019architecture} on experts. In contrast, novices must expend their working memory on new strategies and information to solve problems beyond the capacity of their existing schema on the field \cite{hart2006}. As a result, schema relative to information being processed is a key factor affecting cognitive load. Another factor affecting cognitive load is the medium used to present information \cite{sweller2019architecture}. An example of this factor in practice is how multimedia presentation of information can reduce cognitive load by dividing the loads into different processing segments \cite{sweller2011cognitive}. Using both audio and visual presentation of the same information has been shown to reduce cognitive load by splitting the processing to the visual and auditory segments of human processing \cite{sweller2011cognitive}, instead of overloading a single segment such as visual processing.

\subsection{Cognitive Load Properties}
\subsubsection{Goal-Free Effect}
The goal-free effect states that less cognitive load is received when there’s no specific goal \cite{sweller2019cognitive}. Means-ends analysis is the most generic strategy for solving unfamiliar problems, but it is highly inefficient \cite{hart2006}. Continuously holding in working memory the current state, the goal state, the difference between the two and the legal operations for solving a problem induces a high cognitive load \cite{hart2006}. Without a specific goal in mind, a person is able to perform a similar activity with less cognitive load compared to a person implementing a means-ends strategy. The goal free effect can contribute to learning because cognitive overload could be lessened to allow for schema acquisition.

\subsubsection{Modality Effect}
The modality effect considers the theory that working memory is divided into different sections, such as for visual and auditory \cite{sweller2019cognitive, sweller2011cognitive}. By distributing information into these sections, one is less likely to induce cognitive overload. Cognitive load has been shown to be reduced when presented with the same information when split into audio and visual format compared to only a visual presentation \cite{sweller2011cognitive}.

\subsubsection{Split Attention Effect}
The split-attention effect explains how increased cognitive load is present when a person integrates two sources of information \cite{sweller2019cognitive}. Extraneous cognitive load has been shown to be reduced for those using augmented reality (AR) aids in a physics undergraduate course \cite{klepsch2022}, as information is consolidated realtime using AR technology. 

\subsection{Measuring Cognitive Load}
\subsubsection{Questionnaires}
One methodology in HCI researches is by utilizing questionnaires \cite{hall2024, castanheira2021, ayres2021}. Questionnaires enable subjective responses from the participants, providing HCI studies an overview of the participant and description of experiences and challenges. One famous subjective measure is the NASA-TLX \cite{hall2024, ayres2021} which is famously used in subjective HCI research methodologies to assess mental workload for the above-mentioned aspects. A research instrument such as questionnaires is easily deployable, but it requires extensive use in order to synthesize highly valuable information \cite{castanheira2021}.

\subsubsection{Physiological Measures}
Physiological measures have been shown to be accurate and empirical metrics of cognitive load \cite{lecunff2024, castanheira2021, thees2020}. Eye metrics such as pupil dilation and blink patterns have been shown to be the most sensitive physiological measure to cognitive stimulation \cite{souchet2022, castro-alonso2020modality}. Additionally, other physiological measures of cognitive load exist, such as cerebral \cite{appel2019, babu2018} and cardiovascular \cite{ahmad2023, lecunff2024} measures. These measures provide more insights regarding cognitive load by looking into the physiological reactions of the participants in response to various stimuli and mental tasks. However, compared to questionnaires, physiological measures and instruments are less accessible and difficult to deploy.

\subsection{An Evolving Field of Study}
The cognitive load theory has been applied in various contexts having a common goal: to understand the cognitive effects of the immediate study \cite{sweller2019architecture}. Since its inception, various methods have been implemented and integrated to understand the effects of cognitive load in the performance of human participants in different working tasks \cite{sweller2019architecture}. It is expected that there will be more advancements, and more studies will be conducted that uses the cognitive load theory as their core theoretical framework \cite{sweller2019architecture}.

%%%%%%%%%%%%% SECTION 2 %%%%%%%%%%%%%
\section{Neurodiversity}
Neurodiversity is understood as having different cognitive processes in contrast to neurotypical individuals. It highlights the individual differences in the way that each brain works, and that there is no “correct” way how individuals should think \cite{mack2022, cognitive2025}. In this section, we shall delineate and contextualize neurodiversity in relation to cognitive load, HCI, and computer science.

\subsection{Contextualizing Neurodiversity}
In HCI and computer science, neurodiversity studies frequently look into the subjective aspect, which includes the solicitation of experience and insights from neurodivergent individuals in different settings like school \cite{burtscher2024, cognitive2025} or the workplace \cite{bowman2023, biondi2020, liu2024}. These studies help us understand their unique experiences, preferences, and strengths. With that, this study became more motivated to bridge the gap between empirical methodologies and neurodivergent experiences. Furthermore, the studies also showed us a trend where neurodivergent individuals are often kept from participating in the development of systems and applications that are supposedly dedicated to them \cite{bowman2023, kuang2024, lowy2023, skulmowski2022, spiel2022}, raising a concern in the inclusivity and reliability of the outputs. Consequently, the neurodivergent population, in relation to HCI research, expresses their preference to be part of technological developments \cite{kuang2024, liu2024, skulmowski2022}, providing more valuable insights in the pursuit of a more inclusive technological community. 

\subsection{Neurodiversity and Cognitive Load}
Based on the previous subsections, cognitive load and neurodiversity research, respectively, are plenty. However, studies that delve into the relationship between cognitive load and neurodiversity and their dynamics are highly scarce \cite{vilotijevic2022}, especially how cognitive load plays when neurodiversity is used as an independent variable. A similar study \cite{vilotijevic2022} focusing on the relationship between cognitive load in neurodivergent individuals has discovered and concluded the immense lack of study of those variables through their systematic narrative synthesis. This reinforces the need and impact of a similar study, ensuring an invaluable contribution to the field.

\section{Identified Gaps in Literature and Systems}
There is a lack of study in the area of cognitive load research for neurodivergent individuals  \cite{vilotijevic2022, bowman2023, biondi2020, liu2024}. Research on this area of study could provide opportunities to better serve the population of neurodivergent users in the software industry, as there is a gap in the quality of software for the needs of neurodivergent individuals.

\section{Justification for the Proposal}
The potential of the results of this research will have widespread impact on designing systems for neurodivergent individuals and the general population of software users. The findings of this study could become a basis for designing new conventions and standards for accessible \cite{conrad2023, kosch2023, lowy2023} software design and implementation. Lastly, this research could explore the differences in neurodivergent cognitive processing in relation to the neurotypical, and provide precedence for future research in this largely unexplored area.